% Ege Üniversitesi FBE yüksek lisans tezi — ana dosya.
% Derleme: latexmk main.tex   (ayarlar latexmkrc içinde: pdflatex + BibTeX)
\documentclass[12pt,a4paper]{mythesis}

\usepackage[utf8]{inputenc}
\usepackage[turkish]{babel}
\usepackage{amsmath,amsthm,amsfonts,amssymb}
\usepackage{graphicx}
\usepackage{rotating}
\usepackage{longtable}
\usepackage{lscape}
\usepackage{indentfirst}
\usepackage{fancyhdr}
\usepackage{xkeyval}
\usepackage{hyperref}
\usepackage{natbib}
\usepackage{url}
\usepackage{afterpage}
\usepackage{chngcntr}
\usepackage{xcolor}
\usepackage{tikz}
\usepackage{float}
\usepackage[colorinlistoftodos]{todonotes}
\usepackage[linesnumbered,boxed,lined,ruled,vlined]{algorithm2e}
\usepackage{enumitem}
\usetikzlibrary{positioning,arrows,calc}

\hypersetup{hidelinks}          % bağlantılar tıklanır ama çerçeve görünmez
\graphicspath{{figures/}}       % şekiller kod reposundan PDF olarak buraya gelir
\addto\captionsturkish{\renewcommand{\refname}{KAYNAKLAR DİZİNİ}}

\numberwithin{equation}{section}
\counterwithin{figure}{section}
\setlist{listparindent=\parindent, parsep=0pt}
\def\emptyline{\vspace{11pt}}

\makeatletter
\renewcommand*\l@section{\@dottedtocline{1}{0em}{3em}}
\renewcommand*\l@subsection{\@dottedtocline{2}{0em}{3em}}
\renewcommand*\l@subsubsection{\@dottedtocline{3}{0em}{3em}}
\makeatother

\pagestyle{fancy}
\fancyhead{}
\renewcommand{\headrulewidth}{0.0pt}
\renewcommand{\footrulewidth}{0.0pt}
\chead{\thepage}
\lfoot{}\cfoot{}\rfoot{}
\hyphenpenalty=1000
\tolerance=1000

% Çift taraflı baskı için arkası boş sayfa. Şablon her ön bölümden sonra bunu koyuyor.
% Kılavuz tek taraflıya izin veriyorsa yalnızca bu satır değişir.
\newcommand{\bossayfa}{\newpage\thispagestyle{empty}\mbox{}\newpage}

% ---- Kapak bilgileri (mythesis.cls bunları kapak sayfasına basar) ----
\newcommand{\tezadi}{AZ ETİKETLİ TIBBİ GÖRÜNTÜLERDE TEMEL MODEL DESTEKLİ VERİMLİ ÖĞRENME}
\newcommand{\tezadien}{FOUNDATION MODEL-ASSISTED EFFICIENT LEARNING ON MEDICAL IMAGES WITH LIMITED LABELS}
\title{\textbf{\tezadi}\\[16pt]}
\author{Sefa Samet SÜTÇÜ\\[14pt]}
\advisor{Danışman: Dr. Öğr. Üyesi İsmail IRMAKÇI\\[12pt]}
\kod{Bilgisayar Mühendisliği Anabilim Dalı}
\bolum{Bilgisayar Mühendisliği Yüksek Lisans Programı\\[12pt]}
\location{İzmir \\}
\tarih{2027 \\}

\begin{document}

\begin{titlepage}
\addcontentsline{toc}{section}{İÇ KAPAK}
\pagenumbering{roman}
\end{titlepage}
\setcounter{page}{2}
\maketitle

%%%%%%%%%%%%%%%%% KABUL VE ONAY SAYFASI %%%%%%%%%%%%%%%%%
% Savunmadan sonra doldurulur: tarih, oybirliği/oyçokluğu ve jüri adları.
\thispagestyle{empty}
\vspace*{4cm}

Sefa Samet SÜTÇÜ tarafından yüksek lisans tezi olarak sunulan “\tezadi” başlıklı bu çalışma EÜ Lisansüstü Eğitim ve Öğretim Yönetmeliği ile EÜ Fen Bilimleri Enstitüsü Eğitim ve Öğretim Yönergesi’nin ilgili hükümleri uyarınca tarafımızdan değerlendirilerek savunmaya değer bulunmuş ve ...............\ tarihinde yapılan tez savunma sınavında aday oybirliği/oyçokluğu ile başarılı bulunmuştur.\\

\pagenumbering{roman}
\setcounter{page}{3}
\addcontentsline{toc}{section}{KABUL VE ONAY SAYFASI}
\begin{tabular}{lr}
\textbf{\underline{Jüri üyeleri}} & \hspace{0.5cm}\textbf{\underline{İmza}}\\
  & \\
\textbf{Jüri Başkanı \ \ \ :} Jüri Başkanının Ünvanı-Adı-Soyadı & \\
\\
\textbf{Raportör Üye :} Raportör Üyenin Ünvanı-Adı-Soyadı & \\
\\
\textbf{Üye \ \ \ \ \ \ \ \ \ \ \ \ \ \ : } Üyenin Ünvanı-Adı-Soyadı & \\
\\
\end{tabular}

\bossayfa

%%%%%%%%%%%%%%%%% ETİK KURALLARA UYGUNLUK BEYANI %%%%%%%%%%%%%%%%%
\pagenumbering{gobble}
\clearpage
\thispagestyle{empty}
\pagenumbering{roman}
\setcounter{page}{5}
\addcontentsline{toc}{section}{ETİK KURALLARA UYGUNLUK BEYANI}
\begin{center}\textbf{EGE ÜNİVERSİTESİ FEN BİLİMLERİ ENSTİTÜSÜ}\end{center}
\begin{center}\textbf{ETİK KURALLARA UYGUNLUK BEYANI}\end{center}

\indent EÜ Lisansüstü Eğitim ve Öğretim Yönetmeliğinin ilgili hükümleri uyarınca Yüksek Lisans Tezi olarak sunduğum “\tezadi” başlıklı bu tezin kendi çalışmam olduğunu, sunduğum tüm sonuç, doküman, bilgi ve belgeleri bizzat ve bu tez çalışması kapsamında elde ettiğimi, bu tez çalışmasıyla elde edilmeyen bütün bilgi ve yorumlara atıf yaptığımı ve bunları kaynaklar listesinde usulüne uygun olarak verdiğimi, tez çalışması ve yazımı sırasında patent ve telif haklarını ihlal edici bir davranışımın olmadığını, bu tezin herhangi bir bölümünü bu üniversite veya diğer bir üniversitede başka bir tez çalışması içinde sunmadığımı, bu tezin planlanmasından yazımına kadar bütün safhalarda bilimsel etik kurallarına uygun olarak davrandığımı ve aksinin ortaya çıkması durumunda her türlü yasal sonucu kabul edeceğimi beyan ederim.\\
\\
\\
\indent\indent\indent\indent\indent\indent\indent\indent\indent\indent\indent\quad\text{../../20..}\\ \\
\vspace{.9cm}
\indent\indent\indent\indent\indent\indent\indent\indent\indent\indent\indent\quad\text{İmzası}\\
\vspace{.3cm}
\indent\indent\indent\indent\indent\indent\indent\indent\indent\indent\indent\quad\text{Sefa Samet SÜTÇÜ}

\bossayfa

%%%%%%%%%%%%%%%%% ÖZET %%%%%%%%%%%%%%%%%
\pagenumbering{roman}
\setcounter{page}{7}
\addcontentsline{toc}{section}{ÖZET}
\begin{center}\textbf{ÖZET}\end{center}
\begin{center}\textbf{\tezadi}\end{center}
\begin{center}SÜTÇÜ, Sefa Samet\end{center}
\begin{center}
Yüksek Lisans Tezi, Bilgisayar Mühendisliği Anabilim Dalı\\
Tez Danışmanı: Dr. Öğr. Üyesi İsmail IRMAKÇI\\
Ay Yıl, \pageref{lastpage} sayfa % Yalnızca "Ay Yıl" doldurulur; sayfa sayısı otomatik.
\end{center}

\todo[inline]{Özet, sonuçlar çıktıktan sonra yazılır (Faz 7). En fazla bir sayfa.}

\begin{small}\textbf{Anahtar sözcükler:} temel modeller, bilgi damıtma, kademeli karar, az etiketli öğrenme, kırık tespiti.
\end{small}

\bossayfa

%%%%%%%%%%%%%%%%% ABSTRACT %%%%%%%%%%%%%%%%%
\addcontentsline{toc}{section}{ABSTRACT}
\begin{center}\textbf{ABSTRACT}\end{center}
\begin{center}\textbf{\tezadien}\end{center}
\begin{center}SÜTÇÜ, Sefa Samet\end{center}
\begin{center}
M.Sc. Thesis in Department of Computer Engineering\\
Supervisor: Asst. Prof. Dr. İsmail IRMAKÇI\\
Month Year, \pageref{lastpage} pages
\end{center}

\todo[inline]{Abstract, Özet bittikten sonra çevrilir (Faz 7).}

\begin{small}\textbf{Keywords:} foundation models, knowledge distillation, model cascades, label-efficient learning, fracture detection.
\end{small}

\bossayfa

%%%%%%%%%%%%%%%%% ÖNSÖZ %%%%%%%%%%%%%%%%%
% Şablon notu: tez konusunun belirlenmesinden tamamlanmasına kadar geçen süreç, etken dille.
% Teşekkür bölümü ile karıştırılmamalı.
\addcontentsline{toc}{section}{ÖNSÖZ}\vspace*{1.2cm}
\begin{center}\textbf{\begin{large}ÖNSÖZ\end{large}}\end{center}
\vspace{1.6cm}

\todo[inline]{Önsöz en sonda yazılır.}

\vspace{2cm}
\indent İZMİR\\
\indent ../../20..\\
\\
\indent\indent\indent\indent\indent\indent\indent\indent\textit{Sefa Samet SÜTÇÜ}

\bossayfa


\addcontentsline{toc}{section}{İÇİNDEKİLER}
\tableofcontents
\bossayfa

\addcontentsline{toc}{section}{ŞEKİLLER DİZİNİ}
\begin{center}\textbf{ŞEKİLLER DİZİNİ}\end{center}
\listoffigures
\bossayfa

\addcontentsline{toc}{section}{ÇİZELGELER DİZİNİ}
\listoftables
\bossayfa

%%%%%%%%%%%%%%%%% SİMGELER VE KISALTMALAR DİZİNİ %%%%%%%%%%%%%%%%%
% Tez ilerledikçe alfabetik sırayla güncellenir.
\addcontentsline{toc}{section}{SİMGELER VE KISALTMALAR DİZİNİ}
\begin{center}\textbf{SİMGELER VE KISALTMALAR DİZİNİ}\end{center}
\begin{tabular}{l l}
\underline{Kısaltmalar} & \hspace{1cm}\underline{Açıklama} \\ \\
AUROC & \hspace{1cm} ROC eğrisi altındaki alan (Area Under the ROC Curve)\\ \\
CNN & \hspace{1cm} Evrişimli sinir ağı (Convolutional Neural Network)\\ \\
ECE & \hspace{1cm} Beklenen kalibrasyon hatası (Expected Calibration Error)\\ \\
KD & \hspace{1cm} Bilgi damıtma (Knowledge Distillation)\\ \\
VLM & \hspace{1cm} Görüntü-dil modeli (Vision-Language Model)\\ \\
\end{tabular}

\bossayfa

% ---- Ana metin: çift taraflı kenar boşlukları (iç 5 cm, dış 2 cm) ----
\newgeometry{a4paper, twoside, top=4cm, inner=5cm, outer=2cm, bottom=2cm,
  headheight=15pt, headsep=6mm, footnotesep=6mm, footskip=30mm - 1em}
\newpage
\pagenumbering{arabic}

\section{GİRİŞ}
% Şablon notu: amaç, kapsam ve yöntem; bu anabölümde altbölüm başlığı kullanılmaz.

\todo[inline]{Giriş, Eylül 2027'de yazılır. Atıf testi aşağıda; derlemeden sonra silinebilir.}

Pediatrik el bileği grafileri için kamuya açık en büyük veri seti GRAZPEDWRI-DX'tir \citep{nagyPediatricWristTrauma2022}. Göğüs röntgeniyle öz-denetimli eğitilen RAD-DINO kodlayıcısı \citet{perez-garciaExploringScalableMedical2025} tarafından tanıtılmıştır.

\newpage
\section{LİTERATÜR}
\todo[inline]{Kasım 2026–Şubat 2027: temel modeller, az etiketli öğrenme, bilgi damıtma, kademeli karar ve devretme, kalibrasyon, dağılım kayması, GRAZPEDWRI-DX çalışmaları.}

\newpage
\section{MATERYAL VE YÖNTEM}
\todo[inline]{Her modül bittikçe yazılır: veri ve kohort, öğretmen ve öğrenci modeller, bilgi damıtma, kademeli karar, metrikler, istatistik.}

\newpage
\section{BULGULAR}
\todo[inline]{Mayıs–Eylül 2027: AS1, AS2, AS3. Tablo ve şekiller kod reposundan gelir.}

\newpage
\section{TARTIŞMA}
\todo[inline]{Eylül 2027: bulguların anlamı, sınırlılıklar.}

\newpage
\section{SONUÇ}
% Şablon notu: genel sonuçlar öz, açık ve seçik yazılır.
\todo[inline]{Eylül 2027.}


\newpage
\thispagestyle{empty}\mbox{}
\vspace*{-2.1cm}
\setlength{\headsep}{1.5cm}
\addcontentsline{toc}{section}{KAYNAKLAR DİZİNİ}
\bibliographystyle{bibtez}
\bibliography{references}

%%%%%%%%%%%%%%%%% TEŞEKKÜR %%%%%%%%%%%%%%%%%
\newpage
\addcontentsline{toc}{section}{TEŞEKKÜR}\vspace*{0cm}
\chead{\thepage}
\begin{center}\textbf{\begin{large}TEŞEKKÜR\end{large}}\end{center}

\todo[inline]{Teşekkür en sonda yazılır.}

\indent\indent\indent\indent\indent\indent\indent\indent\indent\indent\indent\quad\text{../../20..}\\ \\
\vspace{.9cm}
\indent\indent\indent\indent\indent\indent\indent\indent\indent\indent\indent\quad\text{İmzası}\\
\vspace{.3cm}
\indent\indent\indent\indent\indent\indent\indent\indent\indent\indent\indent\quad\text{Sefa Samet SÜTÇÜ}
\newpage

%%%%%%%%%%%%%%%%% ÖZGEÇMİŞ %%%%%%%%%%%%%%%%%
% Şablon notu: üçüncü tekil şahısla; vatandaşlık, doğum yeri ve tarihi, yayınlar.
\newpage
\setlength{\headsep}{0.7cm}
\addcontentsline{toc}{section}{ÖZGEÇMİŞ}
\chead{\thepage}
\begin{center}\textbf{\begin{large}ÖZGEÇMİŞ\end{large}}\end{center}

\todo[inline]{Özgeçmiş en sonda yazılır; SIU 2027 bildirisi yayın olarak eklenir.}


\label{lastpage}
\end{document}
